\problema{Lowest Common Ancestor of a Binary Tree}{236}{src/02-arboles/236-lowest-common-ancestor.cpp}{M}

Nos dan un árbol binario y dos nodos \texttt{p} y \texttt{q} que están
garantizados presentes en el árbol. Queremos encontrar su \emph{ancestro
común más bajo} (LCA, \emph{Lowest Common Ancestor}): el nodo más profundo
que tiene tanto a \texttt{p} como a \texttt{q} entre sus descendientes (un
nodo se considera descendiente de sí mismo).

\paragraph{La idea, con un ejemplo de la vida real.} Antes que nada, unas
palabras nuevas. Un \emph{árbol binario} es como un árbol genealógico dibujado
al revés: arriba está el jefe de la familia (lo llamamos \emph{raíz}) y de él
cuelgan personas. Cada persona es un \emph{nodo}, y cada nodo puede tener como
mucho dos hijos: uno a la izquierda y otro a la derecha. Un \emph{ancestro} de
alguien es cualquiera que esté por encima de él en el árbol (su papá, su abuelo,
su bisabuelo\ldots). El \emph{ancestro común más bajo} de dos personas es el
pariente compartido más cercano a ellas, el primero donde se juntan sus dos
líneas familiares.

Imagina que buscas al pariente común de dos primos en ese árbol genealógico.
Te paras en el jefe de la familia y preguntas: «¿alguno de los dos que busco
está en mi lado izquierdo?» y «¿alguno está en mi lado derecho?». Si uno está a
la izquierda y el otro a la derecha, entonces tú, justo donde estás parado, eres
el punto donde se juntan sus caminos: eres la respuesta. Si los dos están del
mismo lado, la respuesta está más abajo por ese lado, así que dejas que esa rama
siga buscando y te avise.

\begin{center}
\ttfamily
\begin{tabular}{ccccccc}
 & & & 3 & & & \\
 & 5 & & & & 1 & \\
6 & & 2 & & 0 & & 8\\
 & 7 & & 4 & & & \\
\end{tabular}
\end{center}

\noindent En este árbol, $LCA(5,1)=3$ porque cada uno cuelga de una rama
distinta de la raíz. En cambio $LCA(5,4)=5$ porque 5 es ancestro directo de
4 (4 vive dentro del propio subárbol de 5). Y $LCA(6,4)=5$ porque tanto 6
como 4 cuelgan del subárbol de 5, aunque en ramas distintas dentro de él.

\paragraph{Cómo lo resuelve el código, paso a paso.} Usamos una función que se
llama a sí misma para revisar cada nodo. Que una función se llame a sí misma se
llama \emph{recursión}: es como un mensajero que, para responder una pregunta,
manda a dos mensajeros más pequeños (uno a la rama izquierda y otro a la derecha)
y espera lo que ellos le contesten. Empezamos por abajo del todo y vamos subiendo
con las respuestas; a ese orden (primero los hijos, luego el padre) se le llama
\emph{post-order}. Cada llamada devuelve «el nodo interesante que encontró»:
\texttt{p}, \texttt{q}, la respuesta ya hallada, o \emph{nada}
(\texttt{nullptr}) si por ahí no había ninguno.
\begin{enumerate}
  \item Si el nodo actual está vacío, o es \texttt{p}, o es \texttt{q}, lo
        devolvemos tal cual y dejamos de bajar por ahí: ya encontramos uno de
        los buscados (o se acabó la rama).
  \item Si no, preguntamos a la rama izquierda y a la derecha.
  \item Si \emph{ambas} ramas devuelven algo, quiere decir que \texttt{p} y
        \texttt{q} quedaron uno a cada lado: el nodo actual es el ancestro común
        más bajo, así que se devuelve a sí mismo.
  \item Si solo una rama devuelve algo, esa es la respuesta y la pasamos hacia
        arriba sin cambiarla.
\end{enumerate}
El código completo (abre el archivo con clic en el título):

\lstinputlisting{src/02-arboles/236-lowest-common-ancestor.cpp}

\complejidad{$O(n)$}{$O(h)$}

\paragraph{¿Qué significa esa complejidad?} Aquí $n$ es la cantidad de nodos y
$h$ es la \emph{altura} del árbol (cuántos pisos tiene desde la raíz hasta el
nodo más hondo). El tiempo $O(n)$ quiere decir que, en el peor caso, miramos
cada nodo una sola vez. La memoria $O(h)$ es porque la recursión va apilando
llamadas mientras baja: nunca hay más llamadas amontonadas que la profundidad
del árbol, como una torre de platos que solo llega tan alto como el árbol es
hondo.

\paragraph{Consejos para la entrevista (Amazon).} Es una pregunta muy común.
Un detalle bonito que puedes mencionar: no comparamos los números de los nodos,
sino su «dirección en memoria» (el puntero), así que funciona aunque haya
valores repetidos. Otra cosa: si el árbol estuviera ordenado (un \emph{árbol de
búsqueda binario}, donde los valores menores van a la izquierda y los mayores a
la derecha), habría un truco aún más rápido que baja directo hacia el lado
correcto sin revisar las dos ramas. No olvides los casos raros: que \texttt{p}
sea ancestro de \texttt{q} (o al revés), o que los dos sean hermanos que
cuelgan del mismo padre.
